\chapter{Maximum Flow}

Picture a network of pipes welded together at junctions. One junction, the
\emph{source}, has water gushing in from outside; another, the \emph{sink}, lets
water drain out. Every pipe has a capacity -- a maximum number of liters per
second it can carry. The water obeys two rules: it never exceeds a pipe's
capacity, and at every junction other than the source and sink, whatever flows in
must flow out. The question this chapter answers, completely and with proof, is:
how much water can you push from source to sink?

That single question -- maximum flow -- turns out to be one of the most leveraged
problems in all of combinatorial optimization. Bipartite matching, edge-disjoint
paths, project selection, image segmentation, baseball elimination, and a long
tail of scheduling problems all reduce to it. And the answer comes wrapped in one
of the cleanest theorems in the subject: the maximum amount you can push equals
the cheapest way to sever the network. Flow and cut are two faces of one coin. We
will earn that statement, not just assert it.

\chapterline

\section{Flow Networks}
\label{ch12:networks}

\begin{definition}[Flow network]
\label{ch12:def-network}
A \emph{flow network} is a directed graph $\Gph=(V,E)$ together with a
nonnegative \emph{capacity} $c(u,v)\ge 0$ for every ordered pair $(u,v)$, two
distinguished vertices the \emph{source} $s$ and the \emph{sink} $t$ with $s\neq
t$. We adopt the convention $c(u,v)=0$ whenever $(u,v)\notin E$. We assume every
vertex lies on some $s$--$t$ path (vertices that do not can be deleted without
changing anything).
\end{definition}

\begin{definition}[Flow]
\label{ch12:def-flow}
A \emph{flow} is a function $f:V\times V\to\R$ satisfying:
\begin{itemize}
  \item \textbf{Capacity:} $f(u,v)\le c(u,v)$ for all $u,v$.
  \item \textbf{Skew symmetry:} $f(u,v)=-f(v,u)$ for all $u,v$.
  \item \textbf{Conservation:} for every $u\in V\setminus\{s,t\}$,
        $\displaystyle\sum_{v\in V} f(u,v)=0$.
\end{itemize}
The \emph{value} of $f$ is the net flow leaving the source,
\[
  \abs{f} \;=\; \sum_{v\in V} f(s,v).
\]
A \emph{maximum flow} is a flow of largest possible value.
\end{definition}

\begin{intuitionbox}
Skew symmetry is a bookkeeping trick, not extra physics. If $3$ units physically
travel from $u$ to $v$, we write $f(u,v)=+3$ and $f(v,u)=-3$. Conservation then
reads ``the signed flows out of every internal vertex sum to zero,'' which is
exactly Kirchhoff's current law. The trick lets one equation handle both
directions of a pipe and makes the residual graph (below) fall out cleanly. When
we draw pictures we only label the positive numbers; the negatives are implicit.
\end{intuitionbox}

A few sanity facts. The empty flow $f\equiv 0$ is always feasible, so a maximum
flow exists whenever the value is bounded -- and it is, because
$\abs{f}\le \sum_v c(s,v)$. With rational capacities the optimum is attained; with
real capacities a maximum still exists by compactness (the feasible flows form a
closed bounded polytope, and $\abs{f}$ is linear). The value, measured at the
source, equals the value measured at the sink:

\begin{lemma}[Value is conserved end to end]
\label{ch12:lem-value-sink}
For any flow $f$, $\;\abs{f}=\sum_{v} f(s,v)=\sum_v f(v,t)$.
\end{lemma}

\begin{proof}
Sum conservation over all $u\notin\{s,t\}$: $\sum_{u\neq s,t}\sum_v f(u,v)=0$.
Add the source row to both sides:
$\sum_{u\neq t}\sum_v f(u,v)=\sum_v f(s,v)=\abs{f}$. Now the left side
double-counts: group the terms $f(u,v)$ by unordered pair $\{u,v\}$ with both
endpoints in $V\setminus\{t\}$; each such pair contributes $f(u,v)+f(v,u)=0$ by
skew symmetry. What survives are exactly the pairs touching $t$, i.e.
$\sum_{u\neq t} f(u,t)=-\sum_u f(t,u)=\sum_v f(v,t)$. Hence
$\abs{f}=\sum_v f(v,t)$.
\end{proof}

\subsection{A small network we will carry through the chapter}

Here is the running example. Six vertices, source $s$, sink $t$, the capacities
printed on the arcs.

\begin{center}
\begin{tikzpicture}[>=Stealth,
  every node/.style={font=\small},
  V/.style={circle,draw=nodeblue,fill=blue!4,thick,minimum size=8mm}]
  \node[V,fill=sourceblue!18] (s) {$s$};
  \node[V] (a) [above right=8mm and 24mm of s] {$a$};
  \node[V] (b) [below right=8mm and 24mm of s] {$b$};
  \node[V] (c) [above right=8mm and 24mm of b] {$c$};
  \node[V] (d) [below right=8mm and 24mm of a] {$d$};
  \node[V,fill=targetgreen!22] (t) [below right=8mm and 24mm of c] {$t$};
  \draw[->,thick] (s) -- node[above left]{$10$} (a);
  \draw[->,thick] (s) -- node[below left]{$10$} (b);
  \draw[->,thick] (a) -- node[left]{$2$} (b);
  \draw[->,thick] (a) -- node[above]{$8$} (d);
  \draw[->,thick] (b) -- node[below]{$9$} (c);
  \draw[->,thick] (c) -- node[left]{$6$} (d);
  \draw[->,thick] (c) -- node[below right]{$10$} (t);
  \draw[->,thick] (d) -- node[above right]{$10$} (t);
\end{tikzpicture}
\end{center}

Eyeballing it: $s$ can emit up to $20$. But $t$ absorbs at most $c(c,t)+c(d,t)=
10+10=20$ as well, so the capacity into $t$ does not bind alone. The real
bottleneck is more subtle, and finding it is the point of the algorithms below.
Hold the picture; we will compute its maximum flow three different ways.

\section{Residual Graphs and Augmenting Paths}
\label{ch12:residual}

The engine behind almost every max-flow algorithm is one idea: given a flow that
is not yet maximum, find a way to push a little more, possibly by \emph{rerouting}
flow you already committed. The data structure that exposes both options at once
is the residual graph.

\begin{definition}[Residual capacity, residual graph]
\label{ch12:def-residual}
Given a flow $f$, the \emph{residual capacity} of an ordered pair is
\[
  c_f(u,v) \;=\; c(u,v) - f(u,v).
\]
The \emph{residual graph} $\Gph_f=(V,E_f)$ has an edge $(u,v)$ exactly when
$c_f(u,v)>0$. An \emph{augmenting path} is a directed $s$--$t$ path in $\Gph_f$.
\end{definition}

\begin{intuitionbox}
Each pipe $(u,v)$ with capacity $c$ carrying $f$ units gives rise to two residual
edges. A \emph{forward} edge $(u,v)$ with leftover $c-f$: room to push more. And
a \emph{backward} edge $(v,u)$ with residual $c_f(v,u)=c(v,u)-f(v,u)=0-(-f)=f$:
the ability to \emph{cancel} up to $f$ units of what is already flowing,
effectively sending it back. Backward edges are why greedy choices are never
fatal -- any commitment can be undone later by routing through its reverse.
\end{intuitionbox}

If $P$ is an augmenting path, its \emph{bottleneck} is
$c_f(P)=\min_{(u,v)\in P} c_f(u,v)>0$. \emph{Augmenting} along $P$ means
\[
  f(u,v)\mathrel{+}= c_f(P),\qquad f(v,u)\mathrel{-}= c_f(P)
  \qquad\text{for each }(u,v)\in P.
\]

Figure~\ref{ch12:fig-residual} shows what the residual graph of the running
network looks like after pushing $8$ units along $s\to a\to d\to t$. The three
arcs on that path each spawn a dashed backward edge of capacity $8$, while the
saturated forward arc $a\to d$ has dropped out (residual $0$).

\begin{figure}[H]
\centering
\begin{tikzpicture}[>=Stealth,
  every node/.style={font=\small},
  V/.style={circle,draw=nodeblue,fill=blue!4,thick,minimum size=8mm},
  fwd/.style={->,thick},
  bwd/.style={->,thick,dashed,structurepurple,bend left=18}]
  \node[V,fill=sourceblue!18] (s) {$s$};
  \node[V] (a) [above right=8mm and 24mm of s] {$a$};
  \node[V] (b) [below right=8mm and 24mm of s] {$b$};
  \node[V] (c) [above right=8mm and 24mm of b] {$c$};
  \node[V] (d) [below right=8mm and 24mm of a] {$d$};
  \node[V,fill=targetgreen!22] (t) [below right=8mm and 24mm of c] {$t$};
  % forward residuals
  \draw[fwd] (s) -- node[above left]{$2$} (a);
  \draw[fwd] (s) -- node[below left]{$10$} (b);
  \draw[fwd] (a) -- node[left]{$2$} (b);
  \draw[fwd] (b) -- node[below]{$9$} (c);
  \draw[fwd] (c) -- node[left]{$6$} (d);
  \draw[fwd] (c) -- node[below right]{$10$} (t);
  \draw[fwd] (d) -- node[above right]{$2$} (t);
  % backward residuals along the augmented path s->a->d->t
  \draw[bwd] (a) to node[below,sloped]{$8$} (s);
  \draw[bwd] (d) to node[above,sloped]{$8$} (a);
  \draw[bwd] (t) to node[above,sloped]{$8$} (d);
\end{tikzpicture}
\caption{Residual graph after pushing $8$ along $s\to a\to d\to t$. Solid arcs are
forward residuals $c_f(u,v)=c(u,v)-f(u,v)$; dashed purple arcs are the backward
residuals (cancellation capacity) created by the push. The forward arc $a\to d$ is
saturated and absent.}
\label{ch12:fig-residual}
\end{figure}

\begin{lemma}[Augmentation preserves feasibility and adds value]
\label{ch12:lem-augment}
If $f$ is a flow and $P$ an augmenting path with bottleneck $\Delta=c_f(P)$, the
function $f'$ produced by augmenting is a flow with $\abs{f'}=\abs{f}+\Delta$.
\end{lemma}

\begin{proof}
Skew symmetry is maintained because we update $(u,v)$ and $(v,u)$ in lockstep.
Capacity: on each $(u,v)\in P$ we raise $f$ by $\Delta\le c_f(u,v)=c(u,v)-f(u,v)$,
so $f'(u,v)\le c(u,v)$; the reverse arc only decreases and stays feasible.
Conservation at an internal $u$ on $P$: the path enters $u$ along some arc
$(x,u)$ and leaves along some arc $(u,y)$. Augmenting raises $f(u,y)$ by $\Delta$
and, since $f(x,u)$ rises by $\Delta$, lowers $f(u,x)$ by $\Delta$. The net
outflow $\sum_v f(u,v)$ therefore changes by $+\Delta$ (from $f(u,y)$) plus
$-\Delta$ (from $f(u,x)$), which is $0$, so conservation still holds. Off the path
nothing changes. Finally $s$ is the path's first vertex with exactly one outgoing
path-arc, so $\abs{f'}=\abs{f}+\Delta$.
\end{proof}

\begin{trapbox}
A backward residual edge does \emph{not} mean ``a pipe that runs the wrong way.''
The physical network is unchanged. The residual graph is a momentary accounting of
the moves available \emph{right now}; it changes every time you augment.
Beginners draw the residual edges once and then forget to recompute them after
augmenting. Always update $c_f$ along the path before searching again.
\end{trapbox}

\section{Ford--Fulkerson and the Max-Flow Min-Cut Theorem}
\label{ch12:ff}

The Ford--Fulkerson method is the whole idea in one sentence: \emph{while an
augmenting path exists, augment along it.} It is a \emph{method}, not yet an
algorithm, because it does not say which path to pick -- and that choice controls
everything about its running time.

\begin{algorithm}[H]
\caption{\textsc{Ford--Fulkerson} (generic augmenting-path method)}
\label{ch12:alg-ff}
\KwIn{network $\Gph=(V,E,c)$, source $s$, sink $t$}
\KwOut{a maximum flow $f$}
$f(u,v)\gets 0$ for all $(u,v)$\;
\While{$\Gph_f$ contains an $s$--$t$ path $P$}{
  $\Delta \gets \min_{(u,v)\in P} c_f(u,v)$\;
  \ForEach{$(u,v)\in P$}{
    $f(u,v)\mathrel{+}=\Delta$\;
    $f(v,u)\mathrel{-}=\Delta$\;
  }
}
\Return $f$\;
\end{algorithm}

To prove this terminates with the optimum we need the dual object: cuts.

\subsection{Cuts}

\begin{definition}[$s$--$t$ cut]
\label{ch12:def-cut}
An \emph{$s$--$t$ cut} is a partition $V=S\cup T$ with $s\in S$, $t\in T$. Its
\emph{capacity} is
\[
  c(S,T)\;=\;\sum_{u\in S}\sum_{v\in T} c(u,v),
\]
counting only arcs crossing from $S$ to $T$. The \emph{net flow across the cut} is
$f(S,T)=\sum_{u\in S}\sum_{v\in T} f(u,v)$.
\end{definition}

\begin{lemma}[Flow across any cut equals the value]
\label{ch12:lem-flow-cut}
For every flow $f$ and every $s$--$t$ cut $(S,T)$, $\;f(S,T)=\abs{f}$.
\end{lemma}

\begin{proof}
Start from $\abs{f}=\sum_v f(s,v)$. For each internal $u\in S$, conservation gives
$\sum_v f(u,v)=0$, so we may add all of them without changing the sum:
$\abs{f}=\sum_{u\in S}\sum_{v\in V} f(u,v)$. Split the inner sum by whether $v\in
S$ or $v\in T$. The $v\in S$ part pairs each $f(u,v)$ with $f(v,u)$, both endpoints
in $S$; by skew symmetry these cancel. What remains is exactly
$\sum_{u\in S}\sum_{v\in T} f(u,v)=f(S,T)$.
\end{proof}

\begin{corollary}[Weak duality]
\label{ch12:cor-weak}
For every flow $f$ and every cut $(S,T)$, $\;\abs{f}\le c(S,T)$.
\end{corollary}

\begin{proof}
$\abs{f}=f(S,T)=\sum_{u\in S,v\in T} f(u,v)\le\sum_{u\in S,v\in T} c(u,v)=c(S,T)$,
using $f\le c$ on each crossing arc.
\end{proof}

So every cut is a certificate of an upper bound on flow, and every flow is a
certificate of a lower bound on the minimum cut. If we ever exhibit a flow and a
cut with $\abs{f}=c(S,T)$, both are simultaneously optimal. The next theorem says
this always happens.

\begin{theorem}[Max-flow min-cut, Ford--Fulkerson 1956]
\label{ch12:thm-mfmc}
Let $f$ be a flow. The following are equivalent:
\begin{enumerate}
  \item $f$ is a maximum flow.
  \item $\Gph_f$ contains no augmenting path.
  \item $\abs{f}=c(S,T)$ for some $s$--$t$ cut $(S,T)$.
\end{enumerate}
In particular the maximum flow value equals the minimum cut capacity.
\end{theorem}

\begin{proof}
\textbf{(1)$\Rightarrow$(2).} If an augmenting path existed,
Lemma~\ref{ch12:lem-augment} would produce a strictly larger flow, contradicting
maximality.

\textbf{(2)$\Rightarrow$(3).} Suppose $\Gph_f$ has no $s$--$t$ path. Let
\[
  S=\set{v: v\text{ is reachable from }s\text{ in }\Gph_f},\qquad T=V\setminus S.
\]
Then $s\in S$ and, since $t$ is unreachable, $t\in T$, so $(S,T)$ is a cut. Take
any $u\in S$, $v\in T$. Because $v$ is not reachable, $(u,v)\notin E_f$, i.e.
$c_f(u,v)=c(u,v)-f(u,v)=0$, so $f(u,v)=c(u,v)$: every crossing arc is saturated.
Therefore $f(S,T)=\sum_{u\in S,v\in T} c(u,v)=c(S,T)$. By
Lemma~\ref{ch12:lem-flow-cut}, $\abs{f}=f(S,T)=c(S,T)$.

\textbf{(3)$\Rightarrow$(1).} If $\abs{f}=c(S,T)$ for some cut, then for every flow
$g$, weak duality gives $\abs{g}\le c(S,T)=\abs{f}$, so $f$ is maximum.
\end{proof}

\begin{whybox}
The proof is the algorithm. When Ford--Fulkerson stops, it stops precisely because
no augmenting path exists -- condition (2). The set $S$ of residual-reachable
vertices is then a minimum cut, computed for free by one graph search from $s$ in
$\Gph_f$. So you never run a separate min-cut routine: max-flow hands you the
certificate. This is the prototype of LP duality made combinatorial.
\end{whybox}

\begin{corollary}[Integrality theorem]
\label{ch12:cor-integral}
If all capacities are integers, there is a maximum flow in which every $f(u,v)$ is
an integer, and Ford--Fulkerson finds one.
\end{corollary}

\begin{proof}
Start from $f\equiv 0$ (integral). Each bottleneck $\Delta=c_f(P)$ is a minimum of
integers, hence a positive integer, and augmentation keeps $f$ integral. The value
rises by at least $1$ each round and is bounded by any cut capacity, so the
process halts at an integral maximum flow.
\end{proof}

This corollary is why flow is such a useful modeling hammer: feed it an integer
program of the right shape and it returns an integer solution, no rounding.

\subsection{Tracing Ford--Fulkerson on the running network}

\begin{labbox}
Capacities as in Section~\ref{ch12:networks}. We pick augmenting paths greedily by
DFS and watch the value climb.

\textbf{Round 1.} Path $s\to a\to d\to t$, residuals $10,8,10$, bottleneck $8$.
Push $8$. Now $f(s,a)=8$, $f(a,d)=8$, $f(d,t)=8$; value $=8$.

\textbf{Round 2.} Path $s\to b\to c\to t$, residuals $10,9,10$, bottleneck $9$.
Push $9$. Value $=17$.

\textbf{Round 3.} Hunt for any augmenting path. Through $c\to t$ we still have
$c_f(c,t)=1$, but reaching $c$ needs $b\to c$, which is saturated
($c_f(b,c)=0$); the alternative into $d$ is $a\to d$, also saturated
($c_f(a,d)=0$). From $s$ we can still reach $a$ ($c_f=2$) and through it $b$ via
$a\to b$ ($c_f=2$), but from $b$ the only forward arc $b\to c$ is dead. No
augmenting path survives.

The residual-reachable set from $s$ is $S=\set{s,a,b}$. Crossing arcs $S\to T$:
$a\to d$ (cap $8$, saturated) and $b\to c$ (cap $9$, saturated). Cut capacity
$=8+9=17=\abs{f}$. Confirmed maximum.
\end{labbox}

The maximum flow value is $17$, and the minimum cut $\set{s,a,b}\mid\set{c,d,t}$
has capacity $17$. Note that neither the source's out-capacity ($20$) nor the
sink's in-capacity ($20$) was the bottleneck -- the constriction sits in the
middle, on the two arcs $a\to d$ and $b\to c$.

\begin{center}
\begin{tikzpicture}[>=Stealth,
  every node/.style={font=\small},
  V/.style={circle,draw=nodeblue,fill=blue!4,thick,minimum size=8mm}]
  \node[V,fill=sourceblue!18] (s) {$s$};
  \node[V] (a) [above right=8mm and 24mm of s] {$a$};
  \node[V] (b) [below right=8mm and 24mm of s] {$b$};
  \node[V] (c) [above right=8mm and 24mm of b] {$c$};
  \node[V] (d) [below right=8mm and 24mm of a] {$d$};
  \node[V,fill=targetgreen!22] (t) [below right=8mm and 24mm of c] {$t$};
  \draw[->,thick] (s) -- node[above left]{$8/10$} (a);
  \draw[->,thick] (s) -- node[below left]{$9/10$} (b);
  \draw[->,thick,gray] (a) -- node[left]{$0/2$} (b);
  \draw[->,very thick,trapred] (a) -- node[above]{$8/8$} (d);
  \draw[->,very thick,trapred] (b) -- node[below]{$9/9$} (c);
  \draw[->,thick,gray] (c) -- node[left]{$0/6$} (d);
  \draw[->,thick] (c) -- node[below right]{$9/10$} (t);
  \draw[->,thick] (d) -- node[above right]{$8/10$} (t);
  \begin{scope}[on background layer]
    \node[draw=structurepurple,thick,dashed,rounded corners,fit=(s)(a)(b),
          inner sep=6pt,label=below:{\color{structurepurple}$S$}]{};
  \end{scope}
\end{tikzpicture}
\end{center}

The two red arcs are the saturated min-cut. The label $f/c$ shows flow over
capacity.

\subsection{Does Ford--Fulkerson even terminate?}

With integer capacities, yes (Corollary~\ref{ch12:cor-integral}): the value rises
by $\ge 1$ per augmentation, so at most $\abs{f^\star}$ rounds occur, where
$f^\star$ is the optimum. That gives a running time of $\Oh(E\cdot\abs{f^\star})$
-- each augmenting path costs $\Oh(E)$ to find by BFS/DFS. This is
\emph{pseudopolynomial}: fine when $\abs{f^\star}$ is small, terrible when
capacities are huge.

\begin{trapbox}
The classic adversarial network: four vertices $s,u,v,t$ with $c(s,u)=c(s,v)=
c(u,t)=c(v,t)=10^9$ and a single middle arc $c(u,v)=1$. If your path chooser keeps
picking $s\to u\to v\to t$ and then $s\to v\to u\to t$ (using the backward edge of
the middle arc), each augmentation moves only $1$ unit, and you take $2\times
10^9$ rounds for a flow of value $2\times 10^9$. With \emph{irrational} capacities
one can even make Ford--Fulkerson run forever and converge to the wrong value. The
fix is to stop choosing paths foolishly.
\end{trapbox}

\begin{center}
\begin{tikzpicture}[>=Stealth,
  every node/.style={font=\small},
  V/.style={circle,draw=nodeblue,fill=blue!4,thick,minimum size=8mm}]
  \node[V,fill=sourceblue!18] (s){$s$};
  \node[V] (u) [above right=6mm and 22mm of s]{$u$};
  \node[V] (v) [below right=6mm and 22mm of s]{$v$};
  \node[V,fill=targetgreen!22] (t) [below right=6mm and 22mm of u]{$t$};
  \draw[->,thick] (s)--node[above left]{$10^9$}(u);
  \draw[->,thick] (s)--node[below left]{$10^9$}(v);
  \draw[->,very thick,trapred] (u)--node[right]{$1$}(v);
  \draw[->,thick] (u)--node[above right]{$10^9$}(t);
  \draw[->,thick] (v)--node[below right]{$10^9$}(t);
\end{tikzpicture}
\end{center}

\begin{historybox}
Lester Ford and Delbert Fulkerson, at the RAND Corporation in 1955--56, were
modeling the Soviet rail network's capacity to move material to Eastern Europe -- a
Cold War interdiction problem. The min-cut was the cheapest set of rail links to
sever. Their theorem unified a scattering of earlier results (Menger's theorem on
disjoint paths, K\H{o}nig's on bipartite matching) under one duality. The bad
instances above were charted later, motivating the search for polynomial rules.
\end{historybox}

\section{Edmonds--Karp: Shortest Augmenting Paths}
\label{ch12:ek}

The cure is almost insultingly simple. \emph{Always augment along a shortest
augmenting path} -- fewest edges, found by BFS in $\Gph_f$. Jack Edmonds and
Richard Karp proved (1972) that this single rule makes the method polynomial,
independent of the capacities.

\begin{algorithm}[H]
\caption{\textsc{Edmonds--Karp}}
\label{ch12:alg-ek}
\KwIn{network $\Gph$, source $s$, sink $t$}
$f\gets 0$\;
\While{BFS in $\Gph_f$ finds a shortest $s$--$t$ path $P$}{
  augment $f$ along $P$ by its bottleneck $c_f(P)$\;
}
\Return $f$\;
\end{algorithm}

The analysis rests on how BFS distances in the residual graph behave as we
augment. Write $\delta_f(v)$ for the number of edges on a shortest $s$--$v$ path in
$\Gph_f$ ($\infty$ if unreachable).

\begin{lemma}[Residual distances never decrease]
\label{ch12:lem-ek-mono}
Under Edmonds--Karp, $\delta_f(v)$ is nondecreasing over time for every $v$.
\end{lemma}

\begin{proof}
Suppose augmenting $f\to f'$ strictly decreased some distance; among all such
vertices pick $v$ with the smallest $\delta_{f'}(v)$. Let $u$ be the predecessor
of $v$ on a shortest $s$--$v$ path in $\Gph_{f'}$, so $\delta_{f'}(v)=
\delta_{f'}(u)+1$ and $(u,v)\in E_{f'}$. By the minimal choice of $v$,
$\delta_{f'}(u)\ge\delta_f(u)$.

Case A: $(u,v)\in E_f$ already. Then $\delta_f(v)\le\delta_f(u)+1\le
\delta_{f'}(u)+1=\delta_{f'}(v)$, contradicting the decrease.

Case B: $(u,v)\notin E_f$ but $(u,v)\in E_{f'}$. The only way a residual edge
appears is that we pushed flow along $(v,u)$, and since Edmonds--Karp augments a
\emph{shortest} path, that path used $(v,u)$, giving $\delta_f(v)=\delta_f(u)-1$.
Then $\delta_f(v)=\delta_f(u)-1\le\delta_{f'}(u)-1=\delta_{f'}(v)-2<
\delta_{f'}(v)$, again contradicting a decrease.
\end{proof}

\begin{theorem}[Edmonds--Karp bound]
\label{ch12:thm-ek}
Edmonds--Karp performs $\Oh(VE)$ augmentations, hence runs in $\Oh(VE^2)$ time.
\end{theorem}

\begin{proof}
Call an edge $(u,v)$ \emph{critical} on an augmentation if it is the bottleneck
(its residual hits $0$ and it leaves $E_f$). When $(u,v)$ is critical it lies on a
shortest path, so $\delta_f(u)+1=\delta_f(v)$. The edge can reappear in $E_f$ only
after we push flow back along $(v,u)$ in some later augmentation, which (shortest
path again) requires $\delta_{f'}(v)+1=\delta_{f'}(u)$ at that time. By
Lemma~\ref{ch12:lem-ek-mono}, distances only grew, so
\[
  \delta_{f'}(u)=\delta_{f'}(v)+1\ge \delta_f(v)+1=\delta_f(u)+2.
\]
Thus between two successive times $(u,v)$ is critical, $\delta(u)$ rises by at
least $2$. Since $\delta(u)\in\set{0,1,\dots,V-1}$ while $u$ is reachable (once
$\delta(u)\ge V$ the vertex is gone for good), each edge is critical $\Oh(V)$
times. Every augmentation makes at least one edge critical, and there are $\Oh(E)$
residual edges, so there are $\Oh(VE)$ augmentations. Each costs $\Oh(E)$ for the
BFS, giving $\Oh(VE^2)$.
\end{proof}

\begin{complexitybox}
\textbf{Ford--Fulkerson (generic):} $\Oh(E\,\abs{f^\star})$, pseudopolynomial, may
not terminate on irrational data. \quad
\textbf{Edmonds--Karp:} $\Oh(VE^2)$, strongly polynomial, capacity-independent.
The win is purely from picking shortest paths -- the same code skeleton, one BFS
instead of an arbitrary search.
\end{complexitybox}

\section{Dinic's Algorithm: Level Graphs and Blocking Flows}
\label{ch12:dinic}

Edmonds--Karp augments one path per BFS. That is wasteful: a single BFS computes
distances to \emph{every} vertex, yet we use it for one path. Dinic's algorithm
(Yefim Dinitz, 1970) reuses each BFS to push a whole \emph{blocking flow} -- many
paths -- before recomputing distances. The result is a strict improvement, and on
the networks arising from matching it is dramatically faster.

\subsection{The level graph}

\begin{definition}[Level graph]
\label{ch12:def-level}
Given $\Gph_f$, run BFS from $s$ to get levels $\ell(v)=\delta_f(v)$. The
\emph{level graph} $\Gph_L$ keeps only the \emph{advancing} residual edges -- those
$(u,v)\in E_f$ with $\ell(v)=\ell(u)+1$.
\end{definition}

Every $s$--$t$ path in $\Gph_L$ has exactly $\ell(t)$ edges, the current shortest
distance. So pushing flow within $\Gph_L$ can only use shortest augmenting paths,
all of the same length. Figure~\ref{ch12:fig-level} draws the level graph of the
running network in Dinic's first phase: the columns are the BFS levels, and only
forward-advancing arcs survive.

\begin{figure}[H]
\centering
\begin{tikzpicture}[>=Stealth,
  every node/.style={font=\small},
  V/.style={circle,draw=nodeblue,fill=blue!4,thick,minimum size=8mm},
  lvl/.style={font=\small\itshape,structurepurple},
  drop/.style={->,thick,gray,dashed}]
  \node[V,fill=sourceblue!18] (s) {$s$};
  \node[V] (a) [above right=6mm and 22mm of s] {$a$};
  \node[V] (b) [below right=6mm and 22mm of s] {$b$};
  \node[V] (c) [above right=6mm and 22mm of b] {$c$};
  \node[V] (d) [below right=6mm and 22mm of a] {$d$};
  \node[V,fill=targetgreen!22] (t) [below right=6mm and 22mm of c] {$t$};
  % advancing edges
  \draw[->,thick] (s) -- node[above left]{$10$} (a);
  \draw[->,thick] (s) -- node[below left]{$10$} (b);
  \draw[->,thick] (a) -- node[above]{$8$} (d);
  \draw[->,thick] (b) -- node[below]{$9$} (c);
  \draw[->,thick] (c) -- node[below right]{$10$} (t);
  \draw[->,thick] (d) -- node[above right]{$10$} (t);
  % dropped level-flat edges
  \draw[drop] (a) -- node[left]{$2$} (b);
  \draw[drop] (c) -- node[left]{$6$} (d);
  % level labels
  \node[lvl] at ($(s)+(0,-1.5)$) {$\ell=0$};
  \node[lvl] at ($(a)+(0,1.0)$) {$\ell=1$};
  \node[lvl] at ($(c)+(0,1.0)$) {$\ell=2$};
  \node[lvl] at ($(t)+(0,-1.5)$) {$\ell=3$};
\end{tikzpicture}
\caption{Level graph of the running network in Dinic phase~1. Solid arcs advance
one level ($\ell(v)=\ell(u)+1$) and stay; the dashed level-flat arcs $a\to b$ and
$c\to d$ (both endpoints at the same level) are excluded. Every $s$--$t$ path here
has length $\ell(t)=3$.}
\label{ch12:fig-level}
\end{figure}

\begin{definition}[Blocking flow]
\label{ch12:def-blocking}
A flow $g$ in the level graph $\Gph_L$ is \emph{blocking} if every $s$--$t$ path in
$\Gph_L$ contains at least one edge saturated by $g$. Equivalently: you cannot
augment any more \emph{within $\Gph_L$}.
\end{definition}

\begin{intuitionbox}
A blocking flow is not a maximum flow of $\Gph_L$ -- it can leave usable capacity
on non-shortest routes. It only promises that no path of the \emph{current
shortest length} survives unsaturated. That is exactly enough: once we saturate
all shortest paths, the next BFS will find that $\delta_f(t)$ has strictly
increased.
\end{intuitionbox}

\subsection{The phase structure}

Dinic proceeds in \emph{phases}. Each phase: (1) BFS to build $\Gph_L$; (2) find a
blocking flow in $\Gph_L$ and add it to $f$. Repeat until $t$ is unreachable.

\begin{lemma}[Distance strictly increases per phase]
\label{ch12:lem-dinic-phase}
After augmenting $f$ by a blocking flow of its level graph, $\delta_f(t)$ strictly
increases.
\end{lemma}

\begin{proof}
Let the phase start with $\delta_f(t)=k$. By monotonicity
(Lemma~\ref{ch12:lem-ek-mono}, whose proof only used that we augmented along
shortest paths -- true of every path in $\Gph_L$), no shortest distance ever
shrinks, so afterward $\delta_{f'}(t)\ge k$. Suppose for contradiction
$\delta_{f'}(t)=k$. A length-$k$ $s$--$t$ path in $\Gph_{f'}$ must climb from level
$0$ to level $k$ in $k$ steps, hence every step goes from level $i$ to level
$i+1$: it is built entirely of advancing edges, so it lives in $\Gph_L$. (New
residual edges created this phase point from level $i$ to $i-1$ and cannot occur on
a length-$k$ ascent.) But the blocking flow saturated some edge on every such
path. Contradiction; hence $\delta_{f'}(t)\ge k+1$.
\end{proof}

\begin{corollary}
\label{ch12:cor-dinic-phases}
Dinic runs at most $V-1$ phases.
\end{corollary}

\begin{proof}
$\delta_f(t)$ starts at $\ge 1$, strictly increases each phase, and once it
exceeds $V-1$ no simple $s$--$t$ path exists, so $t$ is unreachable and we stop.
\end{proof}

\subsection{Finding a blocking flow}

Within a level graph we push flow with a DFS that has one crucial optimization:
the \emph{current-arc} pointer (dead-end pruning). When a DFS from $u$ finds that
an edge leads to a dead end, it never tries that edge again this phase. This is
implemented by keeping, for each $u$, a pointer \texttt{it[u]} into its
advancing-edge list that only ever moves forward.

\begin{algorithm}[H]
\caption{\textsc{Dinic} -- one phase, blocking flow by current-arc DFS}
\label{ch12:alg-dinic}
\KwIn{levels $\ell(\cdot)$, residual $c_f$, pointers \texttt{it[}\,$\cdot$\,\texttt{]}}
\SetKwFunction{Dfs}{Dfs}
\SetKwProg{Fn}{Function}{:}{}
\Fn{\Dfs{$u$, pushed}}{
  \lIf{$u=t$}{\Return pushed}
  \While{\texttt{it[}$u$\texttt{]} points at a valid edge $(u,v)$}{
    \If{$c_f(u,v)>0$ \textbf{and} $\ell(v)=\ell(u)+1$}{
      $d\gets$ \Dfs{$v$, $\min(\text{pushed},\,c_f(u,v))$}\;
      \If{$d>0$}{
        $f(u,v)\mathrel{+}=d$;\quad $f(v,u)\mathrel{-}=d$\;
        \Return $d$ \tcp*{keep \texttt{it[}$u$\texttt{]}: $(u,v)$ may still have room}
      }
    }
    advance \texttt{it[}$u$\texttt{]} \tcp*{$(u,v)$ is dead this phase}
  }
  \Return $0$\;
}
\BlankLine
\While{\Dfs{$s$, $\infty$} $>0$}{ \tcp*{push until the level graph is blocked}}
\end{algorithm}

\begin{whybox}
Why the current-arc pointer makes a phase cheap. Each DFS call from $s$ either
reaches $t$ (a successful push, saturating at least one edge on the path) or hits a
dead end at some vertex $u$, at which point \texttt{it[}$u$\texttt{]} advances past
the offending edge \emph{permanently for this phase}. So every residual edge is
``retreated past'' at most once per phase. The successful paths have length
$\le V$, and at least one edge saturates per successful path, so the total work is
bounded by the edge events: $\Oh(VE)$ per phase.
\end{whybox}

\begin{theorem}[Dinic complexity]
\label{ch12:thm-dinic}
Dinic's algorithm computes a maximum flow in $\Oh(V^2 E)$ time.
\end{theorem}

\begin{proof}
By Corollary~\ref{ch12:cor-dinic-phases} there are $\Oh(V)$ phases. We show each
phase's blocking flow costs $\Oh(VE)$. In the current-arc DFS, charge work to two
kinds of events. \emph{Advances of \texttt{it[}$u$\texttt{]}}: each residual edge
is skipped at most once per phase, total $\Oh(E)$. \emph{Saturations}: each
successful augmenting DFS path has length $\le V$ and saturates at least one of its
edges (the bottleneck), so the number of successful paths is $\Oh(E)$ and the work
along them is $\Oh(VE)$. Summing, one phase is $\Oh(VE)$ and the whole run is
$\Oh(V^2E)$.
\end{proof}

\begin{labbox}
Dinic on the running network. \textbf{Phase 1}: BFS gives levels
$\ell(s)=0$, $\ell(a)=\ell(b)=1$, $\ell(c)=\ell(d)=2$, $\ell(t)=3$. Advancing
edges include $s\to a,s\to b,a\to d,b\to c,c\to t,d\to t$ (the arc $a\to b$ is
level-flat, dropped; $c\to d$ is level-flat, dropped). The blocking-flow DFS pushes
$s\to a\to d\to t$ (bottleneck $8$) and $s\to b\to c\to t$ (bottleneck $9$), then
$s\to a$ has $2$ left but $a$'s only advancing arc $a\to d$ is saturated, dead.
Value after phase $1$ is $17$. \textbf{Phase 2}: BFS in the new residual graph
cannot reach $t$ (the cut $\set{s,a,b}$ is saturated), so $\delta_f(t)=\infty$ and
we stop. Two phases, optimal value $17$ -- compare the three rounds of greedy
Ford--Fulkerson.
\end{labbox}

\begin{historybox}
Dinitz devised this in 1970 for a problem set assigned by Georgy Adelson-Velsky
(of AVL-tree fame), reportedly without having seen Edmonds--Karp; the level-graph
and blocking-flow framing was his. For two decades the algorithm was known in the
West as ``Dinic's'' through a garbled transliteration. The blocking-flow
abstraction later spawned a whole family: Karzanov's $\Oh(V^3)$ blocking flows, the
Goldberg--Tarjan push--relabel view (next chapter), and ultimately the near-linear
interior-point and electrical-flow algorithms of the 2010s.
\end{historybox}

\subsection{Special cases where Dinic is much faster}

\begin{theorem}[Unit-capacity networks]
\label{ch12:thm-unit}
On a network with all capacities $1$, Dinic runs in $\Oh(E\sqrt{E})$ time. On the
unit networks arising from bipartite matching it runs in $\Oh(E\sqrt{V})$.
\end{theorem}

\begin{proof}[Proof sketch]
With unit capacities each blocking flow costs only $\Oh(E)$ (every edge saturates
the first time it carries flow, so the DFS is linear). The phase count is the
subtle part. Run the first $\sqrt{E}$ phases normally; afterward $\delta_f(t)\ge
\sqrt{E}$. The flow still to be pushed, $\abs{f^\star}-\abs{f}$, decomposes into
$s$--$t$ paths in the residual graph, each of length $\ge\sqrt{E}$ and (by unit
capacity) edge-disjoint, so there are at most $E/\sqrt{E}=\sqrt{E}$ of them; each
remaining phase pushes $\ge 1$ unit, bounding the rest to $\sqrt{E}$ phases. Total
$\Oh(\sqrt{E})$ phases $\times\,\Oh(E)$ each. For matching networks
edge-disjointness becomes vertex-disjointness and $\abs{f^\star}\le V$, giving the
$\sqrt V$ bound.
\end{proof}

The $\Oh(E\sqrt V)$ figure is exactly the Hopcroft--Karp bound for bipartite
matching, which is Dinic on the matching network in disguise -- a theme we develop
in the matching chapters.

\section{Capacity Scaling}
\label{ch12:scaling}

There is a different polynomial route that keeps the simple
augment-along-a-path structure but tames large capacities directly: \emph{scaling}.
The idea is to push fat augmentations first.

\begin{intuitionbox}
If you only allow yourself augmenting paths whose bottleneck is at least some
threshold $\Delta$, every augmentation moves $\ge\Delta$ units, so there cannot be
too many of them. Start $\Delta$ at the largest power of two $\le C_{\max}$, run
Ford--Fulkerson restricted to residual edges of capacity $\ge\Delta$, then halve
$\Delta$ and repeat. By the time $\Delta=1$ you are doing ordinary
Ford--Fulkerson, but the flow is already nearly optimal, so few augmentations
remain.
\end{intuitionbox}

\begin{algorithm}[H]
\caption{\textsc{Capacity Scaling}}
\label{ch12:alg-scaling}
\KwIn{network with integer capacities, $C_{\max}=\max_{u,v} c(u,v)$}
$f\gets 0$;\quad $\Delta\gets 2^{\floor{\log_2 C_{\max}}}$\;
\While{$\Delta\ge 1$}{
  \While{$\Gph_f$ has an $s$--$t$ path $P$ using only edges with $c_f\ge\Delta$}{
    augment $f$ along $P$\;
  }
  $\Delta\gets\Delta/2$\;
}
\Return $f$\;
\end{algorithm}

\begin{theorem}[Scaling bound]
\label{ch12:thm-scaling}
Capacity scaling makes $\Oh(E\log C_{\max})$ augmentations and runs in
$\Oh(E^2\log C_{\max})$ time.
\end{theorem}

\begin{proof}
There are $1+\floor{\log_2 C_{\max}}=\Oh(\log C_{\max})$ scaling phases. We bound
the augmentations \emph{within one phase} at threshold $\Delta$ by $\Oh(E)$. At the
end of the $\Delta$-phase, no augmenting path uses only edges with $c_f\ge\Delta$;
let $(S,T)$ be the residual-reachable cut in the graph of such edges. Every
crossing arc has $c_f(u,v)<\Delta$, so the residual capacity across the cut is
$<E\cdot\Delta$. Hence $\abs{f^\star}-\abs{f}<E\Delta$ at that moment. Entering the
next phase, the threshold is $\Delta/2$, and each augmentation there moves
$\ge\Delta/2$ units, so the number of augmentations is at most
$(E\Delta)/(\Delta/2)=2E=\Oh(E)$. Multiplying by $\Oh(\log C_{\max})$ phases gives
$\Oh(E\log C_{\max})$ augmentations, each $\Oh(E)$ to find.
\end{proof}

\begin{complexitybox}
\textbf{Scaling:} $\Oh(E^2\log C_{\max})$. It is \emph{weakly} polynomial -- the
$\log C_{\max}$ factor depends on the numbers -- whereas Edmonds--Karp and Dinic
are \emph{strongly} polynomial. Scaling shines when $E$ is small but capacities
huge, and the technique generalizes beautifully to min-cost flow (later chapters).
\end{complexitybox}

\begin{labbox}
Run scaling on the bad network of the trap above ($c(s,u)=c(u,t)=c(s,v)=c(v,t)=
10^9$, middle arc $c(u,v)=1$). $C_{\max}=10^9$, so $\Delta$ starts at $2^{29}$. The
first phase, with $\Delta=2^{29}$, ignores the middle arc entirely (its capacity
$1<\Delta$) and pushes $2^{29}$ along $s\to u\to t$ and $2^{29}$ along
$s\to v\to t$. Continuing to halve $\Delta$, each phase pushes the two fat disjoint
paths and \emph{never} routes through the $1$-capacity arc until $\Delta=1$, where
a single unit may cross. The whole computation is $\Oh(\log 10^9)\approx 30$
augmentations -- versus $2\times 10^9$ for naive Ford--Fulkerson. That is the
entire point of scaling, visible on one tiny graph.
\end{labbox}

\section{Putting the Bounds Side by Side}
\label{ch12:summary}

\begin{center}
\begin{tabular}{@{}lll@{}}
\toprule
Algorithm & Time & Notes \\
\midrule
Ford--Fulkerson (generic) & $\Oh(E\,\abs{f^\star})$ & pseudopoly; can loop on reals \\
Edmonds--Karp & $\Oh(VE^2)$ & shortest path, strongly poly \\
Capacity scaling & $\Oh(E^2\log C_{\max})$ & weakly poly, fat paths first \\
Dinic & $\Oh(V^2E)$ & blocking flows, $\Oh(V)$ phases \\
Dinic, unit capacity & $\Oh(E\sqrt{E})$ & saturate-once DFS \\
Dinic, bipartite matching & $\Oh(E\sqrt{V})$ & = Hopcroft--Karp \\
\bottomrule
\end{tabular}
\end{center}

For dense graphs Dinic's $\Oh(V^2E)$ beats Edmonds--Karp's $\Oh(VE^2)$; for sparse
huge-capacity graphs scaling wins; for unit/matching networks Dinic is
near-$\sqrt{}$-linear. The next chapter develops push--relabel, which reaches
$\Oh(V^3)$ and $\Oh(V^2\sqrt E)$ by abandoning paths altogether and working with
local ``preflows.''

\begin{connectionbox}
\textbf{Where this plugs into the rest of the book.}
\begin{itemize}
  \item \emph{Min-cut} (the global-minimum-cut chapter): the $s$--$t$ min-cut is a
        free byproduct here via residual reachability; the global min-cut runs $s$--$t$
        cuts over many sinks or uses Stoer--Wagner.
  \item \emph{Matching} (bipartite and general): bipartite matching is integral
        max-flow on a unit network; Dinic gives Hopcroft--Karp's $\Oh(E\sqrt V)$.
  \item \emph{Min-cost flow} and \emph{weighted matching}: replace ``shortest by
        edges'' with ``cheapest by cost,'' and scaling reappears as cost-scaling.
  \item \emph{Laplacians and electrical flows} (later chapters): the modern
        $\Oh(E^{1+o(1)})$ max-flow algorithms route flow by repeatedly solving
        Laplacian systems -- electrical flow is the continuous relaxation of
        combinatorial flow. Max-flow min-cut is the combinatorial shadow of the
        Laplacian's spectral structure, and the residual graph is the discrete
        gradient you will meet again as the $\Lap$ operator.
  \item \emph{LP duality:} Theorem~\ref{ch12:thm-mfmc} is strong duality for the
        flow LP; weak duality (Corollary~\ref{ch12:cor-weak}) is the easy half.
\end{itemize}
\end{connectionbox}

\begin{applicationbox}
\textbf{Vertex capacities and undirected edges} are the two modeling moves you
will reach for constantly. To cap a vertex $v$ at $\kappa$, split it into
$v_{\text{in}}\to v_{\text{out}}$ with capacity $\kappa$ and route all in-arcs to
$v_{\text{in}}$, all out-arcs from $v_{\text{out}}$. To model an undirected edge of
capacity $c$, add both $(u,v)$ and $(v,u)$ with capacity $c$; a max-flow uses at
most one direction net. With these two gadgets and integrality, max-flow solves
edge- and vertex-disjoint paths (Menger), bipartite matching (K\H{o}nig), project
selection / minimum closure, image segmentation (Boykov--Kolmogorov graph cuts),
and baseball elimination. The flow-modeling chapter is a catalog of such
reductions.
\end{applicationbox}

\section{Exercises}
\label{ch12:exercises}

\begin{exercise}[Trace by hand]
On the running network of Section~\ref{ch12:networks}, run Edmonds--Karp,
recording each BFS shortest path, its bottleneck, and the running value. Confirm
the final value is $17$ and that the residual-reachable set from $s$ is
$\set{s,a,b}$. How many augmentations did Edmonds--Karp use, versus the three of
the greedy DFS trace?
\end{exercise}

\begin{exercise}[Min-cut from a max-flow]
Given a maximum flow $f$, prove that the set $S$ of vertices reachable from $s$ in
$\Gph_f$ defines a minimum cut, and that the set of vertices that can \emph{reach}
$t$ in $\Gph_f$ defines (by its complement) another, possibly different, minimum
cut. Construct a network with at least two distinct minimum cuts.
\end{exercise}

\begin{exercise}[Integrality, sharpened]
Show that if every capacity is an even integer, some maximum flow assigns even flow
to every edge. Generalize: if capacities are multiples of $k$, a maximum flow
exists with all flows multiples of $k$. Where in the argument does
Corollary~\ref{ch12:cor-integral} get used?
\end{exercise}

\begin{exercise}[Break Ford--Fulkerson]
Implement the four-vertex trap network with the $1$-capacity middle arc and a
path-chooser that alternates $s\to u\to v\to t$ and $s\to v\to u\to t$. Confirm it
takes $\Th(C_{\max})$ augmentations. Then switch to BFS path selection and show it
finishes in two augmentations. Explain the difference using
Lemma~\ref{ch12:lem-ek-mono}.
\end{exercise}

\begin{exercise}[Edmonds--Karp tightness]
The $\Oh(VE^2)$ bound counts $\Oh(VE)$ augmentations. Construct a family of
networks (e.g.\ layered graphs) where the number of augmentations is $\Om(VE)$,
demonstrating that the analysis is tight up to constants. (Hint: force each edge to
become critical $\Om(V)$ times.)
\end{exercise}

\begin{exercise}[Blocking flow $\neq$ max flow of the level graph]
Give a level graph and a blocking flow $g$ on it that is \emph{not} a maximum flow
of the level graph (some $s$--$t$ capacity remains unused on a longer route).
Explain why Dinic does not care, citing Lemma~\ref{ch12:lem-dinic-phase}.
\end{exercise}

\begin{exercise}[Current-arc pointer correctness]
Prove that resetting the current-arc pointers \texttt{it[}$\cdot$\texttt{]} at the
start of each phase (but never within a phase) yields a correct blocking flow. Then
exhibit a bug: forgetting to reset between phases. What goes wrong, and is the
output still a valid flow (just not maximum), or is it infeasible?
\end{exercise}

\begin{exercise}[Unit-capacity phases]
Fill in the proof of Theorem~\ref{ch12:thm-unit}: after $\sqrt{E}$ phases of Dinic
on a unit network, $\delta_f(t)\ge\sqrt{E}$; conclude the remaining flow value is
$\le\sqrt{E}$ using edge-disjointness of the residual augmenting paths. Where do
you use that capacities are exactly $1$?
\end{exercise}

\begin{exercise}[Scaling correctness]
Prove that capacity scaling terminates with a maximum flow, not merely a large one.
(Hint: the final $\Delta=1$ phase is ordinary Ford--Fulkerson; quote
Theorem~\ref{ch12:thm-mfmc}.) Then re-derive the $\Oh(E\log C_{\max})$ augmentation
bound, being explicit about why the residual cut across the $\Delta$-frontier has
capacity $<E\Delta$.
\end{exercise}

\begin{exercise}[Vertex capacities and Menger]
Using the split-vertex gadget from the application box, model a network where each
\emph{vertex} (not edge) has a throughput limit, and prove the gadget preserves the
max-flow value exactly. Then formulate Menger's vertex-connectivity theorem as a
max-flow statement and prove it as a corollary of Theorem~\ref{ch12:thm-mfmc}.
\end{exercise}

\begin{exercise}[Antisymmetric vs.\ nonnegative flows]
We defined flow with skew symmetry ($f(u,v)=-f(v,u)$). An alternative is
``nonnegative flows'': $f(u,v)\ge 0$ on real arcs with conservation $\sum_v f(v,u)=
\sum_w f(u,w)$. Show the two formulations describe the same feasible values, and
translate the residual-graph construction between them. Which is cleaner for
implementation, and why do most libraries store an edge and its reverse adjacently
in an edge array?
\end{exercise}

\begin{exercise}[Lower bounds on edges]
Extend the model so each edge has a \emph{lower bound} $\ell(u,v)\le f(u,v)\le
c(u,v)$. Describe the standard reduction to ordinary max-flow (a super-source and
super-sink absorbing the excess demands) and state when a feasible flow exists.
Prove that your feasibility condition is both necessary and sufficient.
\end{exercise}
